\centerline{\bf \Huge EZ Соответствия.}

\begin{flushright}
        \scriptsize{Z – это Самость,\\ а Самость – это Z.\\Arc Warden.}
\end{flushright}


\problem{0.1}
Докажите, что натуральное число является точным квадратом тогда и только тогда, когда оно имеет нечётное число делителей.

\problem{0.2}
В выпуклом n-угольнике (n > 3) отметили все точки пересечения диагоналей. Известно, что никакие три диагонали не пересекаются в одной точке. Сколько точек было отмечено?

\problem{0.3}
Среди автобусных билетов с номерами от 000000 до 999999 каких больше: счастливых (у которых сумма первых трёх цифр равна сумме последних трёх цифр) или билетов с суммой цифр 27?

\problem{1}
Каких чисел среди шестизначных больше: тех, у которых каждая цифра больше предыдущей, или тех, у которых каждая цифра меньше предыдущей?

\problem{2}
Каких прямоугольников с целыми сторонами больше: с периметром $2022$ или с периметром $2024$? (Прямоугольники $a \times b$ и $b \times a$ считаются одинаковыми.)

\problem{3}
В группе ЭлМШ учится $49$ боксеров и один борец. Назовём \textit{бандой} любую группу, состоящую из двух или более детей из ЭлМШ. Каких банд больше: с борцом или без борца, и на сколько?

\problem{4}
Вова выписал все последовательности из $2023$ нулей и единиц. Каких последовательностей получилось больше --- тех, в которых чётное число единиц, или тех, в которых количество единиц нечётное?

\problem{5}
Докажите, что число разложений натурального числа $n$ в сумму различных натуральных слагаемых равно числу разложений числа $n$ в сумму нечетных (возможно, повторяющихся) натуральных слагаемых.